\documentclass[conference]{IEEEtran}
\usepackage{times}
\usepackage{graphicx}
% numbers option provides compact numerical references in the text. 
\usepackage[numbers]{natbib}
\usepackage{multicol}
\usepackage[bookmarks=true]{hyperref}
\usepackage{amsmath}
\usepackage{upgreek}

\pdfinfo{
   /Author (Homer Simpson)
   /Title  (Robots: Our new overlords)
   /CreationDate (D:20101201120000)
   /Subject (Robots)
   /Keywords (Robots;Overlords)
}

\begin{document}

% paper title
\title{Template paper for the \\Robotics: Science and Systems Conference}

% You will get a Paper-ID when submitting a pdf file to the conference system
\author{Author Names Omitted for Anonymous Review. Paper-ID [add your ID here]}

%\author{\authorblockN{Michael Shell}
%\authorblockA{School of Electrical and\\Computer Engineering\\
%Georgia Institute of Technology\\
%Atlanta, Georgia 30332--0250\\
%Email: mshell@ece.gatech.edu}
%\and
%\authorblockN{Homer Simpson}
%\authorblockA{Twentieth Century Fox\\
%Springfield, USA\\
%Email: homer@thesimpsons.com}
%\and
%\authorblockN{James Kirk\\ and Montgomery Scott}
%\authorblockA{Starfleet Academy\\
%San Francisco, California 96678-2391\\
%Telephone: (800) 555--1212\\
%Fax: (888) 555--1212}}


% avoiding spaces at the end of the author lines is not a problem with
% conference papers because we don't use \thanks or \IEEEmembership


% for over three affiliations, or if they all won't fit within the width
% of the page, use this alternative format:
% 
%\author{\authorblockN{Michael Shell\authorrefmark{1},
%Homer Simpson\authorrefmark{2},
%James Kirk\authorrefmark{3}, 
%Montgomery Scott\authorrefmark{3} and
%Eldon Tyrell\authorrefmark{4}}
%\authorblockA{\authorrefmark{1}School of Electrical and Computer Engineering\\
%Georgia Institute of Technology,
%Atlanta, Georgia 30332--0250\\ Email: mshell@ece.gatech.edu}
%\authorblockA{\authorrefmark{2}Twentieth Century Fox, Springfield, USA\\
%Email: homer@thesimpsons.com}
%\authorblockA{\authorrefmark{3}Starfleet Academy, San Francisco, California 96678-2391\\
%Telephone: (800) 555--1212, Fax: (888) 555--1212}
%\authorblockA{\authorrefmark{4}Tyrell Inc., 123 Replicant Street, Los Angeles, California 90210--4321}}


\maketitle

\begin{abstract}
The abstract goes here.
\end{abstract}

\IEEEpeerreviewmaketitle

\section{Introduction}
This documents provides a detail explanation of the state-space algorithm that is implemented in this project.
The algorithm is developed based on existing state-space algorithm \cite{Massarini551005}\cite{Massarini921362} and is very similar
to the approach used by PLECS \cite{Allmeling794588} and DSpace \cite{Kiffe9570187}. 



\section{Recognize Circuit Topology}
Although the paper published by PLECS \cite{Allmeling794588} points out how to retrieve state-space matrix from a network topology matrix, 
but the authors did not give any detail on a systematic method for obtaining those state-space matrix. 
In Dspace's paper\cite{Kiffe9570187}, the authors point to a highly influential paper\cite{Kuh1445921} on how to retrieve state-space matrix based on circuit network Topology.
In Massarini and Reggiani's paper, the authors points to the book co-author by L.O.Chua and P.M.Lin\cite{LOChua} for methods on retriving state-space matrix base on circuit network Topology.

After a detail comparison, the author found many of the ideas by Kuh and Rohrer \cite{Kuh1445921} exist in the book \cite{LOChua}. Thus, the author decide
to further develop the algorithm base on the work by Massarini and Reggiani\cite{Massarini551005}, \cite{Massarini921362}.


The paper by Massarini and Reggiani\cite{Massarini921362} start develop based on the \eqref{eq:1} which obtained in \cite{LOChua}.
For a detail overview of the equation, please refer to chapter 4 and chapter 6 of \cite{LOChua}. 
\begin{equation}
  \resizebox{0.5\textwidth}{!}{%
  $
  \begin{bmatrix}
    Day_{11} & 0 & 1 & 0 & D_{ab} & 0\\
    0 & -(D_{zb})^T & 0 & 1 & 0 & -(D_{ab})T\\ 
    \hat{F_{iy} }& \hat{F_{vz}}&\hat{F_{ia}}&\hat{F_{vb}}&\hat{F_{ib}}&\hat{F_{va}}
  \end{bmatrix}
  \times
  \begin{bmatrix}
      i_{y}\\
      V_{z}\\ 
      i_{a}\\ 
      v_{b}\\
      i_{b}\\
      v_{a}
  \end{bmatrix}
  =0
  $%
  }
  \label{eq:1}
  \end{equation}


\section{Extension of the Work}

At the end of section 2in \cite{Massarini921362}, the author mentioned both M and Q is not full rank when there 
exist dependency between state variables. However, the author did not give any specific details in handling those situation while maintaining the number of state variables
consistent through out the simulation. 


Assume given matrix E, M, A, B of a LTI(linear time-invariant) state-space system with the form 
\begin{equation}
  M\times E \times \hat{x} = A \times x + B \times u
  \label{eq:2}
\end{equation} 
where E is the Capacitance/Inductance of the coreesponding element. If M is full rank, it implies the result of \[  ME = M_{f}\] is invertible,
which simplifies the state-space equation in \eqref{eq:2} to 
\begin{equation}
  \hat{x} = M_{f}^{-1} (A\times x + B \times u), M_{f} = ME
  \label{eq:state-space-simplified}
\end{equation}

Consider the following state-space form for an hypothetically RLC circuit. Clearly, M\textsubscript{f} is always invertible regardless the value of L and C.

\resizebox{0.5\textwidth}{!}{%
$
\begin{bmatrix}
  1 & 0 \\
  0 & 1
\end{bmatrix}
\begin{bmatrix}
  L & 0 \\
  0 & C \\
\end{bmatrix}
\hat{x}
=\begin{bmatrix}
  -0.001 & 10 \\
  1 & 2 \\
\end{bmatrix}
x
+
\begin{bmatrix}
  2 \\
  3\\
\end{bmatrix}
u
$%
}


However, the equation in \eqref{eq:state-space-simplified} no longer holds when M is no longer full rank.
In this case, several possible solution is possible.

In \cite{Verghese1102763}, it recognized such system as a "Singular Perturbed Systems." In \cite{Verghese1102763}, the authors proposed filling M matrix with small $\varepsilon$
value . However, the process of determine this $\varepsilon$ is often a iterative, trial-error process.

On the other hand, in \cite{Sajewski}, \cite{Galvao03102018}, and \cite{Li10363677} views such singular system as descriptor state space system. In both \cite{Sajewski} and \cite{Galvao03102018},
it proposed to use either shuffle algorithm, Weierstrass-Kronecker Decomposition or Drazin Inverse method to transform the descriptor state space system into standard state-space form.
However, the transformation involves redefine the state variable as \[ 
\begin{bmatrix}
  z \\
  w\\
\end{bmatrix} = V \times x_{i} \] 
where V is a matrix and x\textsubscript{i} is the independent state in the system. In other words,
the new states is linear combination of independent states of the original system. While this method works well for large system containing parallel capacitors or series inductor, it is not the 
perfect solution for state-space simulation. Because for every switch/diode change, it transformed into a new topoloy with own set of M, A, B matrix, which also changes the definition of state variable.
Given the new topology's initial state value depends on the old topology state value, it would be an issue if the old toplolgy's z defined by different x\textsubscript{i} from the new topology's z.


In this project, I end using a similar simplification that is similiar to shuffle algorithm.
In \eqref{eq:singular-form}, x\textsubscript{1} represents the independent state variable and x\textsubscript{2} represents the dependent state variable.

\begin{equation}
  \begin{bmatrix}
    M_{f11} & M_{f12} \\
    0  & 0  
  \end{bmatrix}
  \begin{bmatrix}
    \hat{x_{1}}\\
    \hat{x_{2}}
  \end{bmatrix}
  = 
  \begin{bmatrix}
    A_{11} & A_{12} \\
    A_{21} & A_{22}    
  \end{bmatrix}
  \begin{bmatrix}
    x_{1} \\
    x_{2}
  \end{bmatrix}
  +
  \begin{bmatrix}
    B1 \\
    B2
  \end{bmatrix}
  u
  \label{eq:singular-form}
\end{equation}


Given A22 is invertable, which is the case when there does not exist CE-loop/I-J cutset \cite{LOChua} in the circuit,
\begin{equation}
  x_{2} = -A_{22}^{-1}A21x_{1}-A_{22}^{-1}B_{2}u
  \label{eq:5}
\end{equation}
and
\begin{equation}
  M_{f11}\hat{x_{1}}+M_{f12}\hat{x_{2}} = A_{11}x_{1}+A_{12}x_{2}+B_{1}u
\end{equation}
which simplifies to
\begin{equation}
  \begin{aligned}
    [M_{f11}-M_{f12}A^{-1}_{22}A_{21}]\hat{x_{1}} = (A_{11}-A_{12}A_{22}^{-1}A_{21})x_{1} \\
     +(B_{1}-A_{12}A_{22}^{-1}B_{2})u
  \end{aligned}
  \label{eq:7}
\end{equation}

//TODO: add exlain why is full rank and invertible

Then, \eqref{eq:7} simplifes to 
\begin{equation}
  \begin{aligned}
  \hat{x_{1}} = [M_{f11}-M_{f12}A^{-1}_{22}A_{21}]^{-1} [ (A_{11}-A_{12}A_{22}^{-1}A_{21})x_{1} \\
  +(B_{1}-A_{12}A_{22}^{-1}B_{2})u ]
  \end{aligned}
  \label{eq:8}
\end{equation} 

In addition, by taking derivative on both side of \eqref{eq:5}, it turn to be
\begin{equation}
  \hat{x_{2}} = -A_{22}^{-1}A21\hat{x_{1}}-A_{22}^{-1}B_{2}\hat{u}
  \label{eq:9}
\end{equation}
Given the simulation has a relative small iteration step, \^{u} is effectively almost 0, which 
simplify the \eqref{eq:9} to
\begin{equation}
  \hat{x_{2}} = -A_{22}^{-1}A21\hat{x_{1}}
\end{equation}



Furthermore, from \cite{Massarini921362}, the author define the output matrix of the system as
\begin{equation}
  Qy=C_{1}\hat{x}+Cx+Du
  \label{eq:11}
\end{equation}
which simplify to 
\begin{equation}
  Qy = C_{1}E^{-1}[Ax+Bu]+Cx+Du
\end{equation}
base on \eqref{eq:2} given M became back to identity matrix  
where Q is an identity matrix when there does not exist CE-loop/I-J cutset.
But when there is dependent in output variables causing Q not being a full rank matrix, we can still simplify the equation to standard form
by finding the minimum norm of the Q matrix




\section{Conclusion} 
\label{sec:conclusion}

The conclusion goes here.

\section*{Acknowledgments}

%% Use plainnat to work nicely with natbib. 

\bibliographystyle{plainnat}
\bibliography{references}

\end{document}


